% ============================================================
% Reusable preamble for LaTeX academic notes / answer-key documents.
% Copy this whole block into a new .tex file, then write \section /
% \subsection content (and \diagramph-style TikZ figures, longtables,
% keybox call-outs) below \begin{document}.
%
% Rename SUBJECT / AUTHOR placeholders below; nothing else needs to change
% to get the same visual style used across this document family.
% ============================================================
\documentclass[11pt,a4paper]{article}

% ---------- Packages ----------
\usepackage[T1]{fontenc}
\usepackage{lmodern}
\usepackage[utf8]{inputenc}
\usepackage[margin=2.3cm]{geometry}
\usepackage{multirow}
\usepackage{array}
\usepackage{booktabs}
\usepackage{longtable}
\usepackage{amsmath}
\usepackage{enumitem}
\usepackage{titlesec}
\usepackage{fancyhdr}
\usepackage{xcolor}
\usepackage{tikz}
\usetikzlibrary{arrows.meta,positioning,shapes.geometric,fit,calc,decorations.pathmorphing}
\usepackage{hyperref}
\usepackage{needspace}
\usepackage{parskip}
\usepackage{tcolorbox}
\tcbuselibrary{skins,breakable}

% ---------- Colors (swap these three for a different palette/subject) ----------
\definecolor{hbrand}{RGB}{31,73,125}   % primary accent -- headers, diagram fills
\definecolor{hgray}{RGB}{90,90,90}     % secondary text -- captions, running heads
\definecolor{hlight}{RGB}{235,241,247} % keybox background tint

% ---------- Section styling ----------
% "Module \thesection" -- change the literal word "Module" to "Unit"/"Chapter"/
% "Week" etc. to match how the source material refers to its own top-level divisions.
\titleformat{\section}{\normalfont\Large\bfseries\color{hbrand}}{Module \thesection}{0.6em}{}
\titlespacing*{\section}{0pt}{18pt}{10pt}
\titleformat{\subsection}{\normalfont\large\bfseries\color{black}}{\thesubsection}{0.6em}{}
\titlespacing*{\subsection}{0pt}{14pt}{6pt}
\titleformat{\subsubsection}{\normalfont\normalsize\bfseries\color{hbrand!80!black}}{}{0em}{}
\titlespacing*{\subsubsection}{0pt}{10pt}{4pt}

% ---------- Header / Footer ----------
\pagestyle{fancy}
\fancyhf{}
\renewcommand{\headrulewidth}{0.4pt}
\lhead{\small\color{hgray}SUBJECT NAME -- DOCUMENT TITLE}  % <-- edit
\rhead{\small\color{hgray}\leftmark}
\cfoot{\thepage}

% ---------- Hyperref / PDF metadata ----------
\hypersetup{
  colorlinks=true,
  linkcolor=hbrand,
  urlcolor=hbrand,
  pdftitle={SUBJECT NAME - DOCUMENT TITLE},   % <-- edit
  pdfauthor={AUTHOR NAME}                     % <-- edit
}

% ---------- Diagram caption wrapper ----------
% Put this directly under every TikZ figure: \diagramcap{Description of the diagram}
\newcommand{\diagramcap}[1]{%
\begin{center}
{\small\itshape\color{hgray} Fig.: #1}
\end{center}
\vspace{2pt}
}

% ---------- Key-point call-out box ----------
% \begin{keybox}[Optional Custom Title]  ...content...  \end{keybox}
\newtcolorbox{keybox}[1][Key Points]{
  colback=hlight, colframe=hbrand!70, boxrule=0.6pt, arc=1.5mm,
  fonttitle=\bfseries\small, coltitle=white, colbacktitle=hbrand,
  title=#1, left=6pt, right=6pt, top=4pt, bottom=4pt, breakable
}

% ---------- TikZ node/arrow styles used by every diagram ----------
% core = the "main"/emphasized node in a diagram (e.g. a source, a master, a core system)
% box  = a standard labeled box
% sbox = a smaller/secondary box (e.g. satellite nodes, table-like cells inside a figure)
% tag  = small italic annotation text next to an arrow or node
% arr / darr = solid / dashed arrow
\tikzset{
  core/.style={draw=hbrand, thick, rounded corners=2pt, fill=hbrand!22, minimum height=0.9cm, align=center, font=\small\bfseries, text=black},
  box/.style={draw=hbrand, thick, rounded corners=2pt, fill=hbrand!10, minimum height=0.9cm, align=center, font=\small, text=black},
  sbox/.style={draw=gray!70, semithick, rounded corners=2pt, fill=gray!8, minimum height=0.75cm, align=center, font=\footnotesize, text=black},
  tag/.style={font=\scriptsize\itshape, text=hgray},
  arr/.style={-{Latex[length=2.2mm,width=1.6mm]}, thick, gray!45!black},
  darr/.style={-{Latex[length=2.2mm,width=1.6mm]}, thick, dashed, gray!45!black}
}

% ---------- List / table conventions ----------
\setlist[itemize]{topsep=2pt,itemsep=2pt,parsep=0pt,leftmargin=1.5em}
\setlist[enumerate]{topsep=2pt,itemsep=2pt,parsep=0pt,leftmargin=1.7em}
\renewcommand{\arraystretch}{1.2}
% Wrapped-text table columns: use L{3cm} / C{3cm} instead of plain l/c in longtable.
\newcolumntype{L}[1]{>{\raggedright\arraybackslash}p{#1}}
\newcolumntype{C}[1]{>{\centering\arraybackslash}p{#1}}

\begin{document}

% ================= TITLE PAGE =================
\begin{titlepage}
\centering
\vspace*{2.3cm}
{\Huge\bfseries\color{hbrand} SUBJECT NAME}\\[0.6cm]           % <-- edit
{\Large DOCUMENT TITLE (e.g. Complete Prep Notes -- Modules 1--3)}\\[1.2cm] % <-- edit
\rule{0.5\linewidth}{0.4pt}\\[1.2cm]
      % <-- edit
{\large \today}\\[2.2cm]
{\large\itshape Topics Covered:}\\[0.4cm]
\begin{minipage}{0.8\linewidth}
\centering
TOPIC 1 \(\cdot\) TOPIC 2 \(\cdot\) TOPIC 3                    % <-- edit
\end{minipage}
\vfill
{\small ONE-PARAGRAPH SCOPE DESCRIPTION -- what this document covers and how it's organized.} % <-- edit
\end{titlepage}

\tableofcontents
\newpage

% Begin content: \section{...} per module/unit, \subsection{...} per lecture/topic/question.

\end{document}
